\documentclass[11pt]{article}

\usepackage{amsmath}
\usepackage{setspace}
\usepackage{amssymb}
\usepackage{changepage}
\usepackage[a4paper, margin=0.5in]{geometry}

\setstretch{1.3}

\hbadness=10000
\setlength{\parindent}{0pt}

\newcommand\numberthis{\addtocounter{equation}{1}\tag{\theequation}}

\begin{document}

    % HEADER 
    \begin{center}
        \LARGE \textbf{STAT3006 Worksheet 01}
    \end{center}

    \rule{\textwidth}{0.3pt}

    % Question 1
    \large \textbf{1. Background: Vector Norms}

    Let $\mathbf{x} \in \mathbb{R}^d$. The $\ell_p$ norm is defined as 

    \begin{align*}
        ||\mathbf{x}||_p \triangleq \left(\sum_{i = 1}^d |x_i|^p\right)^{1/p}
    \end{align*}

    \begin{adjustwidth}{1em}{1em}
    % Question 1a
    (a) Prove that $||\mathbf{x}||_2 \leq ||\mathbf{x}||_1$, for any $\mathbf{x} \in \mathbb{R}^d$
    \\ 

    \textbf{Answer:} For $||\mathbf{x}||_2$ we have \\ 

    \begin{align*}
        ||\mathbf{x}||_2 &= \sqrt{\sum_{i = 1}^d (x_i)^2} \\ 
        ||\mathbf{x}||_2^2 &= \sum_{i = 1}^d (x_i)^2
    \end{align*}

    We have $||\mathbf{x}||$ given by 

    \begin{align*}
        ||\mathbf{x}||_1 &= \sum_{i = 1}^d |x_i| \\ 
        ||\mathbf{x}||_1^2 &= \left(\sum_{i = 1}^d |x_i| \right)^2 \\
                           &= \sum_{i = 1}^d |x_i|^2 + \sum_{i \neq j} |x_i| |x_j| \\ 
                           &\geq ||\mathbf{x}||_2^2
    \end{align*}

    The result follows. \\ 

    % Question 1b
    (b) Infinity norm is defined as $||\mathbf{x}||_\infty \triangleq \max\limits_{i} |x_i|$. Prove
    that 

    \begin{align*}
        ||\mathbf{x}||_\infty = \lim_{p\to\infty}||\mathbf{x}||_p
    \end{align*}

    \textbf{Answer:} Let $M = ||\mathbf{x}||_\infty = \max\limits_{i}|x_i|$. Our goal is to show 
    that 

    \begin{align*}
        \lim_{p\to\infty}\left(\sum_{i = 1}^d |x_i|^p\right)^{1/p} = M
    \end{align*}

    Since $M$ is an upper bound, we know that $|x_i| \leq M, \forall i$. So we have,

    \begin{align*}
        \sum_{i = 1}^d |x_i|^p &\leq \sum_{i = 1}^d M^p \\
                               &= d \cdot M^p \\ 
        \Rightarrow (\sum_{i = 1}^d |x_i|^p)^{1/p} &\leq (d\cdot M^p)^{1/p} \\ 
        \Rightarrow ||\mathbf{x}||_p &\leq d^{1/p} \cdot M \numberthis
    \end{align*}

    Since $M$ is one of the elements in the sum, 

    \begin{align*}
        \sum_{i = 1}^d |x_i|^p &\geq M^p \\ 
        \Rightarrow \left(\sum_{i = 1}^d |x_i|^p\right)^{1/p} &\geq M \\ 
        \Rightarrow ||\mathbf{x}||_p &\geq M
    \end{align*}

    We can now construct an inequality from (1) and (2) and apply Squeeze Theorem

    \begin{align*}
        M \leq ||\mathbf{x}||_p \leq d^{1/p} \cdot M \\ 
        \lim_{p\to\infty} M \leq ||\mathbf{x}||_p \leq \lim_{p\to\infty} d^{1/p} \cdot M \\ 
        M \leq ||\mathbf{x}||_p \leq M
    \end{align*}

    Therefore $||\mathbf{x}||_p = M = ||\mathbf{x}||_\infty$ as $p\to\infty$ \\

    \end{adjustwidth}

    % Question 2
    \textbf{2. Background: Weak and Strong Law of Large Numbers}

    Let $x_1, x_2, \dots$ be IID random numbers with 

    \begin{align*}
        &\mathbb{E}[x_i] = \mu, &\text{Var}(x_i) = \sigma^2 < \infty
    \end{align*}

    and define sample mean 

    \begin{align*}
        \bar{x}_n = \frac{1}{n}\sum_{i = 1}^n x_i
    \end{align*}

    \begin{adjustwidth}{1em}{1em}
    
    % Question 2a
    (a) Show that 

    \begin{align*}
        \mathbb{E}[\bar{x}_n] = \mu \text{ and } \text{Var}(\bar{x}_n) = \frac{\sigma^2}{n} 
    \end{align*}

    \textbf{Answer:} 

    \begin{align*}
        \mathbb{E}[\bar{x}_n] &= \mathbb{E}\left[\frac{1}{n}\sum_{i=1}^n x_i\right] \\ 
                              &= \frac{1}{n}\sum_{i = 1}^n \mathbb{E}[x_i] \\ 
                              &= \frac{n\mu}{n} \\ 
                              &= \mu
                              \\ 
        \text{Var}(\bar{x}_n) &= \text{var}\left(\frac{1}{n}\sum_{i = 1}^n x_i\right) \\ 
                              &= \frac{1}{n^2} \sum_{i = 1}^n \text{Var}(x_i) \\ 
                              &= \frac{n\sigma^2}{n^2} \\ 
                              &= \frac{\sigma^2}{n}
    \end{align*}

    % Question 2b
    (b) \textbf{Convergence in Probability}. If $y$ is a random variable with finite variance,
    then for any $\epsilon > 0$, 

    \begin{align*}
        \mathbb{P}(|y - \mathbb{E}[y]| > \epsilon) \leq \frac{\text{Var}(y)}{\epsilon^2}
    \end{align*}

    This inequality is often referred to as \textit{Chebyshev's Inequality}. Using Chebyshev's 
    Inquality prove that, 

    \begin{align*}
        \mathbb{P}(|\bar{x}_n - \mu| > \epsilon) \to 0
    \end{align*}

    That is, $\bar{x}_n$ converges to $\mu$ in probability, and denoted as 

    \begin{align*}
        \bar{x}_n \xrightarrow{P} \mu
    \end{align*}

    \textbf{Answer:} Substituting results from part (a) into Chebyshev's inequality we have 

    \begin{align*}
        \mathbb{P}(|\bar{x}_n - \mu| > \epsilon) &\leq \frac{\sigma^2}{n \epsilon} \\
        \lim_{n\to\infty}\mathbb{P}(|\bar{x}_n - \mu > \epsilon) &\leq \lim_{n\to\infty}\left(
        \frac{\sigma^2}{n\epsilon}\right) \\ 
                                                                 &= 0
    \end{align*}

    Since $0 \leq \mathbb{P}(\dots) \leq 1$, we have 

    \begin{align*}
        \lim_{n\to\infty}\mathbb{P}(|\bar{x}_n - \mu| > \epsilon) = 0
    \end{align*}

    Therefore $\bar{x}_n$ converges to $\mu$ in probability. \\ 

    % Question 2c
    (c) \textbf{Almost-Sure Convergence}. Another mode of convergence for random variables is that
    of almost-sure convergence, i.e.,
    
    \begin{align*}
        \mathbb{P}\left(\lim_{n\to\infty}\bar{x}_n = \mu\right) = 1
    \end{align*}

    That is, with probability one, the sequence $\bar{x}_n$ converges to $\mu$. \\ 

    Give a brief intuitive explanation of the difference between convergence in probability 
    and almost sure convergence. \\ 

    \textbf{Answer:} With convergence in probably, it says that as $n$ increases, the probability
    that the random variable $\bar{x}_n$ is far from $\mu$ approaches zero. However due to things
    like outliers, even at large $n$ the true probability might not even be exactly zero. Almost
    sure convergence is a stronger condition as it says that once $\bar{x}_n$ reaches $\mu$ at a 
    large-enough $n$, it will stay at $\mu$ with probability 1. \\

    \end{adjustwidth}

    % Question 3 
    \textbf{3. Computation Efficiency and Vectorization}

    We compare two methods for calculating the distance vector $d \in \mathbb{R}^n$ between a 
    query $\mathbf{q} \in \mathbb{R}^d$ and a training set $\{\mathbf{x}_i\}_{i = 1}^n \in 
    \mathbb{R}^d$. Let $\mathbf{X} \in \mathbb{R}^{n x d}$ denote the data matrix whose 
    $i$-th row is $\mathbf{x}_i^T$ \\

    \begin{adjustwidth}{1em}{1em}
    
    % Question 3a
    (a) Recall the algebraic expansion of the squared Euclidean distance to a single point: 

    \begin{align*}
        ||\mathbf{x}_i - \mathbf{q}||_2^2 = \mathbf{x}_i^T \mathbf{x}_i - 2\mathbf{x}_i^T 
        \mathbf{q} + \mathbf{q}^T\mathbf{q} \numberthis
    \end{align*}

    Express the vector of all $n$ squared distances $\mathbf{d}^2 \in \mathbb{R}^n$ using matrix 
    operations such as $\mathbf{X}\mathbf{X}^T$ and $\mathbf{X}\mathbf{q}$. This is often
    called the vectorised approach. \\

    \textbf{Answer:} The vector $\mathbf{d}^2$ can be expressed as 
    
    \begin{align*}
        \mathbf{d}^2 = \begin{bmatrix}
            ||\mathbf{x}_1 - \mathbf{q}||_2^2 \\ \vdots \\ ||\mathbf{x}_n - \mathbf{q}||_2^2
        \end{bmatrix}
    \end{align*}

    From (2) we start with evaluating the first term. It is the squared $\ell_2$ norm of the 
    $i$-th row of $\mathbf{X}$. It can represented by $\mathbf{x}_i^T\mathbf{x}_i = \langle 
    \mathbf{x}_i, \mathbf{x}_i \rangle$. Consider the feature matrix $\mathbf{X}$, we can grab
    every $\langle \mathbf{x}_i, \mathbf{x}_i \rangle, \forall i$ by diag$(\mathbf{XX}^T)$ \\ 

    For the second term, it can trivially be converted to $-2\mathbf{Xq}$. Lastly for the third
    term, it can also be trivially converted to matrix form by using the $n$-dimensional
    identity matrix $\Rightarrow (\mathbf{q}^T\mathbf{q})\mathbf{1}_n$

    \begin{align*}
        \therefore \mathbf{d^2} = \text{diag}(\mathbf{XX^T}) - 2\mathbf{Xq} 
        + \mathbf{q}^T\mathbf{q1}_n
    \end{align*}

    % Question 3b
    (b) The $\textit{naive alternative}$ uses explicit loops over the rows of $\mathbf{X}$ and 
    computes each distance individually:

    \begin{align*}
        &d_i = \sqrt{\sum_{j = 1}^d (\mathbf{X}_{ij} - q_j)^2}, &i = 1,\dots,n
    \end{align*}

    Explain the hardware-level and computational advantages of the vectorised approach over the 
    naive method. \\ 

    \textbf{Answer:} Matrix multiplication involves many independent steps. That is, we don't 
    have to wait
    for the outcome of one row being multiplied to a column to evaluate the current row
    being multiplied to the current column. This allows us to compute the cells of the 
    product in parallel, i.e. SIMD (single instruction, multiple data). Libraries like
    numpy utilize SIMD for matrix multiplication, which is way faster than the naive method
    as it waits for the current row before moving onto the next \\

    \end{adjustwidth}

    % Question 4
    \textbf{4. Programming Benchmarking}

    Implement the two methods discussed in Question 3 and record the time difference for 
    $n = 1,000,000$ and $d = 10$. \\ 

    \textbf{Answer:} Please see code/w1q4.ipynb

\end{document}
